\documentclass[11pt]{article}

\usepackage[T1]{fontenc}
\usepackage[utf8]{inputenc}
\usepackage{amsmath,amssymb,amsthm}
\usepackage{hyperref}
% lstlean.tex — Lean style for the listings package.
% Lineage: Jeremy Avigad (2015), adapted from Assia Mahboubi's SSR style;
% inspired by Olivier Verdier's unixode.sty for the Unicode literate table.
% No CTAN package or upstream repo exists; papers carry this file in their
% source package (cap set arXiv:1907.01449, sphere eversion arXiv:2210.07746).
% Adaptation note: sorry/admit/sorryAx are highlighted in a dedicated error
% tier so that incomplete proofs stand out visually in listings.
\usepackage{listings}
\usepackage{xcolor}

\definecolor{keywordcolor}{rgb}{0.7, 0.1, 0.1}   % red keywords
\definecolor{tacticcolor}{rgb}{0.0, 0.1, 0.6}    % blue tactics
\definecolor{commentcolor}{rgb}{0.4, 0.4, 0.4}   % grey comments
\definecolor{symbolcolor}{rgb}{0.0, 0.1, 0.6}    % blue symbols
\definecolor{sortcolor}{rgb}{0.1, 0.5, 0.1}      % green sorts
\definecolor{errorcolor}{rgb}{0.6, 0, 0}         % dark red: sorry/admit
\definecolor{stringcolor}{rgb}{0.5, 0.1, 0.5}    % purple strings

\lstdefinelanguage{lean}{
  % Anything between $ becomes LaTeX math mode
  mathescape=false,
  % Comments may or not include Latex commands
  texcl=false,
  % keywords, list taken from lean-syntax.txt
  morekeywords=[1]{import, prelude, protected, private, noncomputable,
    definition, meta, renaming, hiding, parameter, parameters, begin,
    constant, constants, lemma, variable, variables, theory, print,
    theorem, example, open, section, namespace, universe, universes, end,
    modifier, modifiers, using, export, exposing, axiom, inductive,
    coinductive, with, structure, class, instance, attribute, local, where,
    by, have, show, let, in, if, then, else, match, do, for, fun, assume,
    from, take, calc, include, omit, infix, infixl, infixr, notation,
    macro, syntax, elab, set_option, initialize, deriving, opaque, abbrev,
    mutual, partial, scoped, termination_by, decreasing_by},
  % Sorts
  morekeywords=[2]{Sort, Type, Prop},
  % tactics, list taken from lean-syntax.txt (tactic keywords)
  morekeywords=[3]{intro, intros, apply, exact, refine, rfl, simp, simp_all,
    simpa, simp_rw, rw, rewrite, erewrite, ring, ring_nf, omega, decide,
    norm_num, norm_cast, induction, cases, rcases, obtain, constructor,
    ext, funext, conv, congr, field_simp, linarith, nlinarith, positivity,
    tauto, trivial, aesop, use, by_cases, by_contra, push_neg, contrapose,
    symm, trans, assumption, contradiction, exfalso, specialize,
    generalize, revert, clear, rename_i, subst, unfold, delta, change,
    convert, split_ifs, interval_cases, fin_cases, zify, gcongr,
    nth_rewrite, suffices, generalizing, at, only, native_decide},
  % warnings - these three tokens mark an unfinished proof; they must never appear at a code position in a delivered file
  morekeywords=[4]{sorry, admit, sorryAx},
  literate=
    {α}{{\ensuremath{\alpha}}}1
    {β}{{\ensuremath{\beta}}}1
    {γ}{{\ensuremath{\gamma}}}1
    {δ}{{\ensuremath{\delta}}}1
    {ε}{{\ensuremath{\varepsilon}}}1
    {ζ}{{\ensuremath{\zeta}}}1
    {η}{{\ensuremath{\eta}}}1
    {θ}{{\ensuremath{\theta}}}1
    {ι}{{\ensuremath{\iota}}}1
    {κ}{{\ensuremath{\kappa}}}1
    {λ}{{\ensuremath{\lambda}}}1
    {μ}{{\ensuremath{\mu}}}1
    {ν}{{\ensuremath{\nu}}}1
    {ξ}{{\ensuremath{\xi}}}1
    {π}{{\ensuremath{\pi}}}1
    {ρ}{{\ensuremath{\rho}}}1
    {σ}{{\ensuremath{\sigma}}}1
    {τ}{{\ensuremath{\tau}}}1
    {υ}{{\ensuremath{\upsilon}}}1
    {φ}{{\ensuremath{\varphi}}}1
    {χ}{{\ensuremath{\chi}}}1
    {ψ}{{\ensuremath{\psi}}}1
    {ω}{{\ensuremath{\omega}}}1
    {Γ}{{\ensuremath{\Gamma}}}1
    {Δ}{{\ensuremath{\Delta}}}1
    {Θ}{{\ensuremath{\Theta}}}1
    {Λ}{{\ensuremath{\Lambda}}}1
    {Ξ}{{\ensuremath{\Xi}}}1
    {Π}{{\ensuremath{\Pi}}}1
    {Σ}{{\ensuremath{\Sigma}}}1
    {Φ}{{\ensuremath{\Phi}}}1
    {Ψ}{{\ensuremath{\Psi}}}1
    {Ω}{{\ensuremath{\Omega}}}1
    {ℵ}{{\ensuremath{\aleph}}}1
    {∀}{{\ensuremath{\forall}}}1
    {∃!}{{\ensuremath{\exists}!}}1
    {∃}{{\ensuremath{\exists}}}1
    {∈}{{\ensuremath{\in}}}1
    {∉}{{\ensuremath{\notin}}}1
    {∘}{{\ensuremath{\circ}}}1
    {∩}{{\ensuremath{\cap}}}1
    {∪}{{\ensuremath{\cup}}}1
    {⊆}{{\ensuremath{\subseteq}}}1
    {⊂}{{\ensuremath{\subset}}}1
    {⊇}{{\ensuremath{\supseteq}}}1
    {⊃}{{\ensuremath{\supset}}}1
    {⊄}{{\ensuremath{\nsubseteq}}}1
    {∅}{{\ensuremath{\emptyset}}}1
    {∧}{{\ensuremath{\wedge}}}1
    {∨}{{\ensuremath{\vee}}}1
    {¬}{{\ensuremath{\neg}}}1
    {⊤}{{\ensuremath{\top}}}1
    {⊥}{{\ensuremath{\bot}}}1
    {↔}{{\ensuremath{\leftrightarrow}}}1
    {→}{{\ensuremath{\rightarrow}}}1
    {←}{{\ensuremath{\leftarrow}}}1
    {↑}{{\ensuremath{\uparrow}}}1
    {⇒}{{\ensuremath{\Rightarrow}}}1
    {⇐}{{\ensuremath{\Leftarrow}}}1
    {⟹}{{\ensuremath{\Longrightarrow}}}1
    {↦}{{\ensuremath{\mapsto}}}1
    {≤}{{\ensuremath{\leq}}}1
    {≥}{{\ensuremath{\geq}}}1
    {≠}{{\ensuremath{\neq}}}1
    {≈}{{\ensuremath{\approx}}}1
    {≡}{{\ensuremath{\equiv}}}1
    {≃}{{\ensuremath{\simeq}}}1
    {≅}{{\ensuremath{\cong}}}1
    {×}{{\ensuremath{\times}}}1
    {·}{{\ensuremath{\cdot}}}1
    {±}{{\ensuremath{\pm}}}1
    {⊕}{{\ensuremath{\oplus}}}1
    {⊗}{{\ensuremath{\otimes}}}1
    {⋆}{{\ensuremath{\star}}}1
    {▸}{{\ensuremath{\blacktriangleright}}}1
    {⟨}{{\ensuremath{\langle}}}1
    {⟩}{{\ensuremath{\rangle}}}1
    {⦃}{{\ensuremath{\{\!\mid}}}1
    {⦄}{{\ensuremath{\mid\!\}}}}1
    {⟦}{{\ensuremath{[\![}}}1
    {⟧}{{\ensuremath{]\!]}}}1
    {⌊}{{\ensuremath{\lfloor}}}1
    {⌋}{{\ensuremath{\rfloor}}}1
    {⌈}{{\ensuremath{\lceil}}}1
    {⌉}{{\ensuremath{\rceil}}}1
    {√}{{\ensuremath{\surd}}}1
    {∣}{{\ensuremath{\mid}}}1
    {∤}{{\ensuremath{\nmid}}}1
    {∎}{{\ensuremath{\blacksquare}}}1
    {⋯}{{\ensuremath{\cdots}}}1
    {…}{{\ensuremath{\ldots}}}1
    {∑}{{\ensuremath{\sum}}}1
    {∏}{{\ensuremath{\prod}}}1
    {⁻¹}{{\ensuremath{^{-1}}}}1
    {¹}{{\ensuremath{^1}}}1
    {²}{{\ensuremath{^2}}}1
    {³}{{\ensuremath{^3}}}1
    {ⁿ}{{\ensuremath{^n}}}1
    {ᵐ}{{\ensuremath{^m}}}1
    {₀}{{\ensuremath{_0}}}1
    {₁}{{\ensuremath{_1}}}1
    {₂}{{\ensuremath{_2}}}1
    {₃}{{\ensuremath{_3}}}1
    {₄}{{\ensuremath{_4}}}1
    {₅}{{\ensuremath{_5}}}1
    {ₙ}{{\ensuremath{_n}}}1
    {ₘ}{{\ensuremath{_m}}}1
    {ₖ}{{\ensuremath{_k}}}1
    {ᵢ}{{\ensuremath{_i}}}1
    {ⱼ}{{\ensuremath{_j}}}1
    {ℤ}{{\ensuremath{\mathbb{Z}}}}1
    {ℕ}{{\ensuremath{\mathbb{N}}}}1
    {ℚ}{{\ensuremath{\mathbb{Q}}}}1
    {ℝ}{{\ensuremath{\mathbb{R}}}}1
    {ℂ}{{\ensuremath{\mathbb{C}}}}1
    {𝔽}{{\ensuremath{\mathbb{F}}}}1
    {𝒫}{{\ensuremath{\mathcal{P}}}}1
    {∗}{{\ensuremath{\ast}}}1
    {•}{{\ensuremath{\bullet}}}1
    {‼}{{!!}}1,
  % Comments
  morecomment=[l]{--},
  morecomment=[s]{/-}{-/},
  morestring=[b]",
  stringstyle=\color{stringcolor},
  keepspaces=true,
  keywordstyle=[1]\color{keywordcolor}\bfseries,
  keywordstyle=[2]\color{sortcolor}\bfseries,
  keywordstyle=[3]\color{tacticcolor}\bfseries,
  keywordstyle=[4]\color{errorcolor}\bfseries\underline,
  escapeinside={(*@}{@*)},
  columns=fullflexible,
  basicstyle=\ttfamily\small,
  commentstyle=\color{commentcolor}\itshape,
  breaklines=true,
  showstringspaces=false,
  frame=single,
  framesep=2mm,
  xleftmargin=4mm,
  numbers=left,
  numberstyle=\tiny\color{commentcolor},
  numbersep=6pt,
}

% Inline Lean code in prose: \lean{Int.ModEq.pow}
\newcommand{\lean}[1]{\lstinline[language=lean,basicstyle=\ttfamily\normalsize]|#1|}

\lstset{language=lean}

% XeTeX rendering patch: listings' literate table does not apply to multibyte
% characters under XeTeX, so the Unicode symbols used in Lean sources are mapped
% to math glyphs via active characters (ASCII-only source, no extra packages);
% xeCJK+SimSun covers any CJK characters inside listings. Rendering-only: the
% listing source bytes are unchanged; the \ifdefined guard keeps the pdflatex
% path unaffected.
\ifdefined\XeTeXversion
  {\lccode`\~="2115 \lowercase{\global\catcode`~=\active \global\def~{\ensuremath{\mathbb{N}}}}}
  {\lccode`\~="2192 \lowercase{\global\catcode`~=\active \global\def~{\ensuremath{\rightarrow}}}}
  {\lccode`\~="2194 \lowercase{\global\catcode`~=\active \global\def~{\ensuremath{\leftrightarrow}}}}
  {\lccode`\~="2200 \lowercase{\global\catcode`~=\active \global\def~{\ensuremath{\forall}}}}
  {\lccode`\~="2227 \lowercase{\global\catcode`~=\active \global\def~{\ensuremath{\wedge}}}}
  {\lccode`\~="2228 \lowercase{\global\catcode`~=\active \global\def~{\ensuremath{\vee}}}}
  {\lccode`\~="203C \lowercase{\global\catcode`~=\active \global\def~{!!}}}
\fi
\ifdefined\XeTeXversion
  \usepackage{xeCJK}
  \setCJKmainfont{SimSun}
\fi

\newtheorem{theorem}{Theorem}

\title{The Diophantine equation $x^2 = n!! \cdot \operatorname{lcm}(1,\dots,n)$: the complete solution set, formalized and machine-checked in Lean 4}
\author{Aurora0134\thanks{Contact: via the GitHub account Aurora0134.}}
\date{\today}

\begin{document}
\maketitle

\begin{abstract}
We determine the complete solution set of $x^2 = n!! \cdot \operatorname{lcm}(1,\dots,n)$ over
the natural numbers: the product is a perfect square if and only if
$n \in \{0,1,2,5,7\}$, with square roots $x = 1,1,2,30,210$. This is the second of the
two equations that Novakovi'c (arXiv:2601.16757) states in his own text as open, and the
answer given here is strictly stronger than the finiteness he asks for. The result, an
iff-characterization together with its parity criterion and witness layer, is formalized in
Lean 4 (12 theorems, 333 lines) and machine-checked: every proof compiles with no
\texttt{sorry}, and \texttt{\#print axioms} reports only
\{propext, Quot.sound, Classical.choice\}. An exclusivity search across OEIS value columns,
arXiv metadata, seven neighbouring full texts, the public Lean ecosystem and mathlib found
no occupying record for the statement itself. This note is a priority claim of record, not a
full paper.
\end{abstract}

\section{Introduction}
We give the complete solution set of the Diophantine equation
$x^2 = n!! \cdot \operatorname{lcm}(1,\dots,n)$, where $n!!$ is the double factorial and
$\operatorname{lcm}(1,\dots,n)$ is the least common multiple of the integers from $1$ to $n$.
The equation is the second of the pair that Novakovi'c (arXiv:2601.16757) lists as open in
his own text; a companion result settled the first equation
$x^2 = n! \cdot n\#$. This note is a priority claim of record for the present result: the
narrative paper follows separately. The formalization was checked by the Lean 4 kernel, with
no \texttt{sorry} and an axiom set limited to the three standard proof-relevant axioms.

\section{Statement}
% LEAN: tasks/20261006-dfact-lcm-square :: DFactLcmSquare.dfact_lcm_square_solution_set  grade=完全证明
\begin{theorem}\label{thm:solutionset}
For every natural number $n$, $n!! \cdot \operatorname{lcm}(1,\dots,n)$ is a perfect square
if and only if $n = 0$ or $n = 1$ or $n = 2$ or $n = 5$ or $n = 7$.
\end{theorem}

% LEAN: tasks/20261006-dfact-lcm-square :: DFactLcmSquare.isSquare_dfact_mul_lcmUpto_iff  grade=完全证明
\begin{theorem}\label{thm:parity}
For every natural number $n$, $n!! \cdot \operatorname{lcm}(1,\dots,n)$ is a perfect square
if and only if, for every prime $p$, the total exponent
$\nu_p(n!!) + \lfloor \log_p n \rfloor$ is even.
\end{theorem}

% LEAN: tasks/20261006-dfact-lcm-square :: DFactLcmSquare.dfact_lcm_witnesses  grade=完全证明
\begin{theorem}\label{thm:witnesses}
The five solution points hold:
$1^2 = 0!! \cdot \operatorname{lcm}(1,\dots,0)$, $1^2 = 1!! \cdot \operatorname{lcm}(1,\dots,1)$,
$2^2 = 2!! \cdot \operatorname{lcm}(1,\dots,2)$, $30^2 = 5!! \cdot \operatorname{lcm}(1,\dots,5)$
and $210^2 = 7!! \cdot \operatorname{lcm}(1,\dots,7)$.
\end{theorem}

\begin{lstlisting}[caption={Formal statements (Lean 4, frozen snapshot \texttt{9baa672f}).},label={lst:stmt}]
theorem isSquare_dfact_mul_lcmUpto_iff (n : ℕ) :
    IsSquare (n‼ * Nat.lcmUpto n) ↔
      ∀ p : ℕ, p.Prime → Even ((n‼).factorization p + p.log n) := by
  rw [Nat.isSquare_iff_even_factorization]
  refine forall_congr' fun p => imp_congr_right fun hp => ?_

theorem dfact_lcm_witnesses :
    1 ^ 2 = 0‼ * Nat.lcmUpto 0 ∧ 1 ^ 2 = 1‼ * Nat.lcmUpto 1 ∧
      2 ^ 2 = 2‼ * Nat.lcmUpto 2 ∧ 30 ^ 2 = 5‼ * Nat.lcmUpto 5 ∧
        210 ^ 2 = 7‼ * Nat.lcmUpto 7 := by
  decide

theorem dfact_lcm_square_solution_set (n : ℕ) :
    IsSquare (n‼ * Nat.lcmUpto n) ↔ n = 0 ∨ n = 1 ∨ n = 2 ∨ n = 5 ∨ n = 7 := by
\end{lstlisting}

\section{Proof}
The complete machine proof is in \texttt{proofs/final-dfact-lcm-square.lean} of this
package (333 lines, 12 theorems). The key excerpt below shows the opening of the main
theorem: the forward direction is by contradiction, splits at $n = 4$, and dispatches the
two remaining ranges to the even arm, the odd arm and the witness layer.

\begin{lstlisting}[caption={Excerpt: the main theorem, forward direction (frozen artifact, sha256 \texttt{1107cad4}).},label={lst:proof}]
theorem dfact_lcm_square_solution_set (n : ℕ) :
    IsSquare (n‼ * Nat.lcmUpto n) ↔ n = 0 ∨ n = 1 ∨ n = 2 ∨ n = 5 ∨ n = 7 := by
  constructor
  · intro h
    by_contra hne
    push_neg at hne
    rcases lt_or_ge n 4 with hlt | hge
\end{lstlisting}

The proof has three parts. Theorem~\ref{thm:parity} reduces squareness to a parity
condition on prime exponents, using $\operatorname{lcm}(1,\dots,n)$'s factorization
$\nu_p(\operatorname{lcm}(1,\dots,n)) = \lfloor \log_p n \rfloor$. Even $n \geq 4$ are
excluded by Bertrand's postulate: a prime $p \in (n/2, n]$ contributes $\nu_p(n!!) = 0$ and
$\lfloor \log_p n \rfloor = 1$, an odd total. Odd $n \geq 33$ are excluded by an interval
prime lemma for the interval $(n/4, n/3]$, proved in the accompanying
library development from a finite
certificate chain and a tail estimate, not as an instance of a published interval
theorem, and imported by the final file rather than restated in it: the prime $p$ found there contributes $\nu_p(n!!) = 2$ and
$\lfloor \log_p n \rfloor = 1$, again an odd total. The remaining band $9 \leq n \leq 32$ and
the small cases $n = 3$ are handled by six explicit non-square witnesses with named
odd-exponent primes. Theorem~\ref{thm:witnesses} supplies the five solution points, which
gives the converse direction.

\section{Verification evidence}
\begin{itemize}
\item Toolchain: Lean 4 v4.34.0, mathlib4 v4.34.0 (rev \texttt{5ed29652}).
\item Statement integrity: the frozen statement file has sha256
      \texttt{9baa672fa40e74626369e4803d66bedd9b808860975a832447f89a2021f0dca0};
      the final proof file has sha256
      \texttt{1107cad44affa0bc84bb4a285c0381d6e6a0afd15c08d935cbda0cee7a256d8e}.
\item Kernel check: strict-mode compilation of the final file exits 0 with no
      \texttt{sorry} (the only diagnostic is one deprecation warning on \texttt{push\_neg}).
      \texttt{\#print axioms} output for the twelve theorems:
\begin{verbatim}
'DFactLcmSquare.isSquare_dfact_mul_lcmUpto_iff' depends on axioms: [propext, Classical.choice, Quot.sound]
'DFactLcmSquare.dfact_lcm_witnesses' depends on axioms: [propext, Classical.choice, Quot.sound]
'DFactLcmSquare.dfact_lcm_square_solution_set' depends on axioms: [propext, Classical.choice, Quot.sound]
'DFactLcmSquare.exists_sq_dfact_lcm_iff' depends on axioms: [propext, Classical.choice, Quot.sound]
\end{verbatim}
\item Audit trail: the frozen statement's kernel well-definedness record, the strict
      compilation log, the axiom print, the statement diff (clean, declaration by
      declaration), the degenerate-probe outputs (the statement is load-bearing: perturbing
      the solution set turns the proof red) and the independent final audit's verdict are in
      \texttt{audit/} of this package.
\item The deposited \texttt{.lean} files are copies in which comments carrying
      pre-publication workflow narrative were blanked; the blanking never touches code, and
      per-file byte equality of the comment-free code view is asserted by the generating
      script (record: \texttt{audit/comment-elision-record.md}).
\end{itemize}

\section{Novelty evidence}
The exclusivity search's four-state adjudication is \textbf{no occupying record found},
with the statement layer clean at every channel that reaches it.

\begin{itemize}
\item \textbf{OEIS value columns.} The solution-$n$ column, the solution-root column
      $x = 1,1,2,30,210$, the object value column
      $a(n) = n!! \cdot \operatorname{lcm}(1,\dots,n)$ (four prefix variants, plus
      alternating-parity sub-columns and drop-prefix cuts in the review pass), the square
      value column and the squarefree-part columns were searched, together with keyword
      layers. The object and solution-root columns are \textbf{null} across all prefix
      variants. The residual hits are unrelated items, or A061742 = $(p_n\#)^2$, which is
      entered precisely because $x = n\#$ and is not a registration of this equation.
\item \textbf{arXiv metadata.} Fourteen load-bearing query forms were run, with the
      positive control \texttt{all:"double factorial"} = 34 hits confirming the channel was
      live: no same form at the metadata layer.
\item \textbf{Neighbouring full texts.} Seven items were read in full (2601.16757,
      2510.24312, 2602.09687, 2602.23838, 2604.09730, 2610.01899, 2606.23485), together
      with all 17 references of the anchor: zero coverage of this equation.
\item \textbf{Citation axis.} Semantic Scholar and OpenAlex, two independent channels, on
      the anchor: \texttt{cited\_by\_count} = 0 on both, i.e.\ the anchor's open case has
      not been closed by anyone citing it.
\item \textbf{Public Lean ecosystem.} The whole public-conjectures repository at content
      level: \texttt{doubleFactorial} occurs in exactly one file (A001818, a different
      object); the statement layer has zero registration.
\item \textbf{mathlib.} Ripgrep with the positive control \texttt{Nat.Prime} = 21 files:
      \texttt{doubleFactorial} occurs in 4 files, none in a number-theoretic square
      context; \texttt{lcmUpto} in one; the two objects never co-occur.
\end{itemize}

Two boundaries are recorded explicitly rather than left implicit. First, the
$\operatorname{lcm}(1,\dots,n)$ side of the object is occupied at the sequence level:
OEIS A056839 registers the squarefree part of $\operatorname{lcm}(1,\dots,n)$. That entry
does not mention the double factorial anywhere, and it registers a sequence, not the
characterizing solution set of this equation. Second, the review layer records that
Crossref is a relevance-ranked channel and that the OpenAlex title search missed the anchor
by indexing status; neither is used as a load-bearing zero hit. The full query record,
including every zero hit, is in \texttt{audit/novelty-record.md} of this package.

\section{Scope and limitations}
The graded conclusion for this result is \emph{complete proof}: the twelve theorems
compile strictly with zero \texttt{sorry} and an axiom set limited to
\{propext, Quot.sound, Classical.choice\}, and the statements match the frozen snapshot
declaration by declaration. Three boundaries are stated rather than implied. The interval
prime lemma for $(n/4, n/3]$ used in the odd arm is proved in the accompanying
library development as a finite
certificate chain plus an Erd\H{o}s-style tail argument; this note does not claim it as an
instance of any published interval theorem, and it is imported rather than shipped, so a
kernel compile of the deposited final file needs those library modules beside it (their
identity is recorded in the audit record of this package). The observation that the solution roots are
$x = n\#$ is a numerical observation recorded alongside the theorem, not a separate
formalized statement. This note contains no literature survey and makes no claim about
related equations: the companion equation $x^2 = n! \cdot n\#$ was settled separately, and
the intersection $\{5,7\}$ of this solution set with the Brocard solutions $\{4,5,7\}$ is a
numerical coincidence with no implication in either direction.

\section*{Data availability}
The Lean sources, the compile and axiom records and the exclusivity-search record are
packaged on Zenodo (DOI: \texttt{10.5281/zenodo.23309733}). The note is released under CC BY 4.0 and
the code under MIT.

\section*{Use of generative AI}
The formalization and proof reported in this note were developed with AI assistance
(large language model: \texttt{kimi-k3}) under human direction and review. All statements
and proofs reported here were checked by the Lean 4 kernel (Lean v4.34.0, mathlib
v4.34.0): every proof compiles with no \texttt{sorry}, and \texttt{\#print axioms} reports
only \{propext, Quot.sound, Classical.choice\}. The author reviewed all content and takes
full responsibility for it.

\section*{Author contributions}
The author selected the problem, verified the semantics of the formal statements against
the informal ones, and reviewed the kernel evidence and the exclusivity record before
signing off this deposit.

\begin{thebibliography}{9}
\bibitem{novakovic2026}
S.~Novakovi'c.
\newblock The Diophantine equation $P(x)=\overset{r}{\underset{i=1}{\prod}}H_{n_i}$.
\newblock arXiv:2601.16757, 2026.

\bibitem{novakovic2025}
S.~Novakovi'c.
\newblock Diophantine equations involving double factorials.
\newblock arXiv:2510.24312, 2025.

\bibitem{novakovic2026b}
S.~Novakovi'c.
\newblock On a generalization of the Brocard--Ramanujan Diophantine equation.
\newblock arXiv:2602.09687, 2026.

\bibitem{novakovic2026c}
S.~Novakovi'c.
\newblock Products of factorials which are products of factorials.
\newblock arXiv:2602.23838, 2026.

\bibitem{novakovic2026d}
S.~Novakovi'c.
\newblock A comment on the equation $n!!=a_1!!\cdots a_t!!$.
\newblock arXiv:2604.09730, 2026.

\bibitem{yudin2026}
F.~Yudin.
\newblock Square products of factorials and a conjecture of Erd\H{o}s and Graham.
\newblock arXiv:2610.01899, 2026.

\bibitem{peixoto2026}
L.~V. Peixoto.
\newblock The Brocard Problem: Structural Invariants, $p$-Adic Density, and the Generative Sieve.
\newblock arXiv:2606.23485, 2026.

\bibitem{lean4}
L.~de Moura and S.~Ullrich.
\newblock The Lean 4 theorem prover and programming language.
\newblock \emph{CADE-28}, 2021. \url{https://doi.org/10.1007/978-3-030-79876-5_37}

\bibitem{mathlib}
The mathlib Community.
\newblock The Lean mathematical library.
\newblock \emph{CPP 2020}, 2020. \url{https://doi.org/10.1145/3372885.3373824}
\end{thebibliography}

\end{document}
